\section{Certified direct claims and closed lower bounds}\label{s:bounds}

\subsection{Verification protocol}\label{ss:protocol}

The computational source of truth contains exactly thirteen direct Rabung
claims: nine three-colour claims for $17\le k\le25$ and four two-colour
claims for $25\le k\le28$.  For each triple $(p,r,k)$ the release protocol
requires:
\begin{enumerate}
\item two executions of the structured triplet \{mirror scan, full scan,
  CPU walk\}, with internal profile equality and equality between executions;
\item two acceptances by the separate full CPU criterion;
\item acceptance by the stand-alone \texttt{verify\_claim}; and
\item for the two primes above the historical Montgomery-safety threshold, an
  additional full-interval witness that uses no Montgomery arithmetic.
\end{enumerate}
The parser is fail-closed: it checks the process status, the exact $p,r,k$,
all profile equalities, (a), (B$^*$), and the final verdict, and rejects any
output containing \texttt{DIFF}, \texttt{FAIL}, \texttt{ECHEC}, or a negative
verdict.  Its digest engages the claim identity, profiles, criterion values,
and final verdict.  The enclosing JSON manifest separately records the source
and binary hashes, compiler version, exact commands, and hashes of every raw
output.  A SHA-256 of that manifest therefore binds the computational claim to
its build and evidence.  The final Markdown and \LaTeX{} claim tables are
generated only from an exact set of thirteen accepted manifests.

\IfFileExists{generated_validated_claims.tex}{%
  \input{generated_validated_claims}%
}{%
  \begin{center}
  \fbox{\begin{minipage}{0.91\linewidth}\small
  \textbf{Validated-claim table unavailable in this build.}
  The canonical table appears only when the exact set of thirteen fail-closed
  claim manifests is available and validated.  The publication contract
  requires the generated table in a release build.
  \end{minipage}}
  \end{center}
}

This rule is stricter than the historical campaign transcript.  The latter
recorded a max-run checksum that did not engage $p$, and some pasted
\texttt{verify\_claim} outputs were truncated.  Those files are retained as
historical evidence, not used as the final source of truth.  The thirteen
claims are documented in \texttt{audit/claims.json}; the expensive final
point-verification run must populate the complete raw manifests before the
release is frozen.  The numerical bounds in this section use only the
thirteen triples, each of which is also re-derived by the separate-source
checker described in Section~\ref{ss:certificates}.

\subsection{Three colours: direct and recurrence-derived bounds}\label{ss:r3}

The historical scan produced nine candidate direct Rabung certificates.  The headline value
\[
  W(3,17)>31\,385\,622\,833
\]
comes from $p=1\,961\,601\,427$ and improves Monroe's published
$15\,509\,557\,937$ by a factor $2.02$~\cite{Monroe2026}.  A direct certificate
is not automatically a general record: published bounds in other colour
columns may imply a stronger result through a recurrence.

Let $q_k$ denote the largest prime not exceeding $k$.  From
Blankenship--Cummings--Taranchuk~\cite[Theorem~2.1]{bct2018},
\[
  W(3,k)>q_k\bigl(W(2,k)-1\bigr).
\]
If the registered input is the strict integer bound $W(2,k)>B$, then
$W(2,k)-1\ge B$, so the generated consequence is exactly
\begin{equation}\label{eq:bct-r3}
  W(3,k)>q_k B,
\end{equation}
with no off-by-one ambiguity.  The baseline program iterates this rule and
Xu's recurrence to a fixed point, retains a provenance path for every generated
candidate, and selects the best bound only afterwards.  Its output also embeds
the source catalogue and SHA-256 hashes of both input registries.

% Generated by tools/close_baselines.py; do not edit by hand.
\begin{table}[ht]
\centering\footnotesize
\begin{tabular}{c r r r l}
\toprule
$k$ & direct $p$ & direct bound & best closed bound & method \\
\midrule
17 & $1\,961\,601\,427$ & $31\,385\,622\,833$ & $31\,385\,622\,833$ & direct Rabung \\
18 & $1\,999\,994\,851$ & $33\,999\,912\,468$ & $33\,999\,912\,468$ & direct Rabung \\
19 & $1\,999\,999\,927$ & $35\,999\,998\,687$ & $35\,999\,998\,687$ & direct Rabung \\
20 & $1\,999\,999\,927$ & $37\,999\,998\,614$ & $37\,999\,998\,614$ & direct Rabung \\
21 & $1\,999\,999\,927$ & $39\,999\,998\,541$ & $39\,999\,998\,541$ & direct Rabung \\
22 & $1\,999\,999\,927$ & $41\,999\,998\,468$ & $49\,058\,715\,046$ & BCT, $q=19$ (published $W(2,k)$) \\
23 & $1\,999\,999\,927$ & $43\,999\,998\,395$ & $142\,741\,265\,701$ & BCT, $q=23$ (published $W(2,k)$) \\
24 & $1\,999\,999\,927$ & $45\,999\,998\,322$ & $252\,542\,143\,876$ & BCT, $q=23$ (published $W(2,k)$) \\
25 & $1\,999\,999\,927$ & $47\,999\,998\,249$ & $628\,673\,328\,287$ & BCT, $q=23$ (archived input) \\
26 & -- & -- & $1\,221\,912\,599\,848$ & BCT, $q=23$ (archived input) \\
27 & -- & -- & $2\,068\,976\,009\,617$ & BCT, $q=23$ (archived input) \\
28 & -- & -- & $2\,159\,051\,058\,266$ & BCT, $q=23$ (archived input) \\
\bottomrule
\end{tabular}
\caption{Direct three-colour Rabung certificates and the best general lower bounds
obtained after closing the registered baselines under the published
Blankenship--Cummings--Taranchuk and Xu recurrences.  A dash means that the
GPU campaign made no direct three-colour claim for that length.  For
$k=22,23,24,25$ the direct certificates are the strongest direct Rabung bounds
located in the dated corpus but are not general records. Every entry is generated with a
machine-readable provenance path in \texttt{audit/generated/bounds\_closure.json}.}
\label{tab:r3}
\end{table}


Thus the direct certificates for $17\le k\le21$ remain the strongest bounds
located in the registered published-and-archived source corpus after recurrence
closure.  For $22\le k\le25$ the direct rows remain improvements within the
direct Rabung family but are not best general bounds.  In particular, the
published two-colour values already dominate the direct rows at $k=22,23,24$.
For $k=25$, even Monroe's published
$W(2,25)>23\,003\,662\,489$ would give
$W(3,25)>529\,084\,237\,247$; the archived input accepted by historical independent checks,
$W(2,25)>27\,333\,622\,969$ strengthens this to the first line below.

% Generated by tools/close_baselines.py; do not edit by hand.
\begin{align*}
W(3,25) &> 628\,673\,328\,287 \\
W(3,26) &> 1\,221\,912\,599\,848 \\
W(3,27) &> 2\,068\,976\,009\,617 \\
W(3,28) &> 2\,159\,051\,058\,266
\end{align*}


The four displayed inequalities are theorem consequences, not additional GPU
claims.  Priority for the two-colour primes remains Monroe's; we report the
candidate re-verification and apply the published recurrence.  The final
reproducibility claim remains conditional on completion of the raw-manifest
release protocol above.  Xu's recurrence is included in the automatic closure
but does not improve any of these three-colour winners.

\subsection{Two colours and archived provenance}\label{ss:r2}

After the scan, we located a table from a second, unpublished phase of Monroe's
BOINC project (2021, two colours only)~\cite{MonroePhase2Archive}.  The author's project page states that
this phase pushed the checked range from one billion to four
billion~\cite{MonroePage}; the largest prime actually listed in the recovered
table is $3\,478\,685\,143$.  We credit the listed primes to Monroe and use our
checker only to establish that the four triples below satisfy Rabung's
criterion.

\begin{table}[ht]
\centering\small
\begin{tabular}{c r r l}
\toprule
$k$ & prime $p$ & $W(2,k)>$ & method and provenance \\
\midrule
25 & $1\,138\,900\,957$ & $27\,333\,622\,969$ & archive and historical recheck; final manifest pending \\
26 & $2\,125\,065\,391$ & $53\,126\,634\,776$ & archived Monroe prime; final manifest pending \\
27 & $3\,459\,826\,103$ & $89\,955\,478\,679$ & historical no-Montgomery recheck; final manifest pending \\
28 & $3\,476\,732\,783$ & $93\,871\,785\,142$ & historical no-Montgomery recheck; final manifest pending \\
\bottomrule
\end{tabular}
\caption{Two-colour succinct Rabung certificates used as recurrence inputs.
The method column separates priority for the archived prime from the candidate
re-verification; final raw manifests are still pending.}
\label{tab:r2}
\end{table}

For $k=25$ the historical campaign summaries identify the same last candidate
valid prime in $[9.7\times10^8,2\times10^9)$ as the archived table; complete
campaign recovery and ordered-prime identity remain release gates.  For
$k=26,27,28$ our campaign-range candidates are weaker; the displayed archive
primes were checked directly in historical runs.
Those runs supply candidate evidence for the listed triples; final release
manifests must make the checks independently auditable.  The computation
does not establish that Monroe's historical phase tested every prime up to
four billion: that coverage statement is an attributed archival claim, not a
consequence of the recovered table.

The source is the Internet Archive capture of 28~April~2022 of
\url{http://vdwnumbers.org/bounds.txt}~\cite{MonroePhase2Archive}; its own header reads ``Last updated:
18~Aug~2021''.  The release records the immutable capture URL, both dates, the
source-repository blob identity, and its SHA-256 in
\texttt{publication/baseline-provenance.json}; it does not redistribute the
third-party table absent permission.  Once the thirteen fail-closed manifests
and their raw outputs are populated, the stand-alone verifier will establish a
separately auditable verification event.  The defensible novelty formulation
for that release is therefore: an independently reproduced verification
accompanied by a stand-alone public checker and complete reproducible logs, not
first discovery of the primes.
